\begin{tikzpicture}[font=\figurefont\small,
  box/.style={draw=BookRule,fill=BookPaper,align=center,
    text width=30mm,minimum height=14mm},
  arr/.style={-{Stealth[length=2mm]},draw=BookInk,line width=.8pt},
  control/.style={draw=BookAmber,dashed,-{Stealth[length=2mm]}}]
  \node[box,fill=FigureModel!12] (model) at (0,0) {学习模型\\$\widehat P,\widehat R$};
  \node[box,fill=FigureOutput!12] (policy) at (5,0) {策略改进\\改变访问分布};
  \node[box,fill=FigureInput!12] (env) at (10,0) {真实执行与校验\\获得新经验};
  \node[box,fill=FigureLoss!12] (fit) at (5,-2.8) {在新访问区域\\重新拟合模型};
  \draw[arr] (model) -- (policy);
  \draw[arr] (policy) -- (env);
  \draw[arr] (env.south) |- (fit.east);
  \draw[arr] (fit.west) -| (model.south);
  \node[align=center,text width=29mm] (penalty) at (0,2.3) {不确定性惩罚\\改变模型评分};
  \node[align=center,text width=29mm] (short) at (5,2.3) {缩短想象时域\\限制传播深度};
  \node[align=center,text width=29mm] (cal) at (10,2.3) {真实数据校准\\需要额外交互};
  \draw[control] (penalty) -- (model);
  \draw[control] (short) -- (policy);
  \draw[control] (cal) -- (env);
  \node[align=center,text width=128mm] at (5,-4.3)
    {离线条件下右侧的真实交互边不可凭空补上。\\固定策略上的低误差，不能自动控制新策略的误差。};
\end{tikzpicture}
